\documentclass[10pt,a4paper]{article}
\usepackage[a4paper,margin=14mm,top=12mm,bottom=14mm]{geometry}
\usepackage{fontspec}
\setmainfont{Liberation Sans}
\usepackage{polyglossia}
\setdefaultlanguage{french}
\usepackage{amsmath}
\usepackage{xcolor}
\usepackage{graphicx}
\usepackage[most]{tcolorbox}
\usepackage{tabularx,booktabs,array,multirow}
\usepackage{enumitem}
\usepackage{fancyhdr}
\usepackage{lastpage}
\usepackage[hidelinks]{hyperref}
\definecolor{uimm}{HTML}{005677}
\definecolor{cyan}{HTML}{13A2DD}
\definecolor{fond}{HTML}{E6F2FA}
\definecolor{ok}{HTML}{2E7D32}
\definecolor{ko}{HTML}{B3261E}
\definecolor{violet}{HTML}{7B1FA2}
\setlength{\parindent}{0pt}
\setlength{\parskip}{3pt}
\setlist{nosep}
\setlist[itemize]{leftmargin=13pt,label=\textcolor{cyan}{\textbullet}}
\newcommand{\titre}[1]{\vspace{5pt}{\color{uimm}\large\bfseries #1}\par\vspace{-3pt}{\color{cyan}\rule{\linewidth}{1.2pt}}\par\vspace{0pt}}
\tcbset{boite/.style={colback=fond,colframe=uimm,boxrule=0.6pt,arc=2pt,left=6pt,right=6pt,top=4pt,bottom=4pt},
        defi/.style={colback=violet!6,colframe=violet,boxrule=0.8pt,arc=2pt,left=6pt,right=6pt,top=4pt,bottom=4pt,
                     title={\bfseries Défi},coltitle=white,colbacktitle=violet,fonttitle=\small}}
\newcommand{\bandeau}[3]{%
\begin{tcolorbox}[colback=uimm,colframe=uimm,arc=2pt,left=4pt,right=8pt,top=4pt,bottom=4pt]
\begin{minipage}[c]{18mm}\setlength{\fboxsep}{1.2pt}\fcolorbox{white}{white}{\includegraphics[height=16.5mm]{logo_uimm.jpg}}\end{minipage}\hspace{4mm}%
\begin{minipage}[c]{\dimexpr\linewidth-22mm\relax}\color{white}
{\small #1}\\[1pt]{\LARGE\bfseries #2}\\[2pt]{\large #3}\hfill{\small CQPM Technicien qualité}
\end{minipage}
\end{tcolorbox}}
\newcommand{\reponse}[1][1]{\par\vspace{2pt}{\color{gray!60}\foreach \i in {1,...,#1}{\rule{\linewidth}{0.3pt}\par\vspace{9pt}}}}
\usepackage{pgffor}
\pagestyle{fancy}\fancyhf{}
\renewcommand{\headrulewidth}{0pt}
\newcommand{\pied}[1]{\fancyfoot[L]{\footnotesize\color{uimm}#1}\fancyfoot[R]{\footnotesize\color{uimm}Page \thepage{} sur \pageref{LastPage}}}

\usepackage{tikz}
\usepackage{longtable}
\renewcommand{\bandeau}[3]{%
\begin{tcolorbox}[colback=uimm,colframe=uimm,arc=2pt,left=4pt,right=8pt,top=4pt,bottom=4pt]
\begin{minipage}[c]{18mm}\setlength{\fboxsep}{1.2pt}\fcolorbox{white}{white}{\includegraphics[height=16.5mm]{logo_uimm.jpg}}\end{minipage}\hspace{4mm}%
\begin{minipage}[c]{\dimexpr\linewidth-22mm\relax}\color{white}
{\small #1}\\[1pt]{\LARGE\bfseries #2}\\[2pt]{\large #3}
\end{minipage}
\end{tcolorbox}}
% en-tête d'une fiche : code, titre, question, bloc, durée, organisation, niveau
\newcommand{\entete}[7]{%
\begin{tcolorbox}[colback=uimm,colframe=uimm,arc=2pt,left=4pt,right=6pt,top=4pt,bottom=4pt]
\begin{minipage}[c]{18mm}\setlength{\fboxsep}{1.2pt}\fcolorbox{white}{white}{\includegraphics[height=16.5mm]{logo_uimm.jpg}}\end{minipage}\hspace{3mm}%
\begin{minipage}[c]{\dimexpr\linewidth-58mm\relax}\color{white}\raggedright
{\small Base SPC · #4}\\[1pt]{\Large\bfseries #1 · #2}\\[2pt]{\itshape #3}
\end{minipage}\hfill
\begin{minipage}[c]{33mm}\color{white}\small\raggedleft
Durée : \textbf{#5}\\ \textbf{#6}\\ Niveau #7 sur 3
\end{minipage}
\end{tcolorbox}}
\newcommand{\tagE}{{\setlength{\fboxsep}{1.5pt}\colorbox{cyan}{\color{white}\scriptsize\bfseries Ellistat}}}
\newcommand{\tagW}{{\setlength{\fboxsep}{1.5pt}\colorbox{uimm}{\color{white}\scriptsize\bfseries Page web}}}
\tcbset{pred/.style={colback=orange!6,colframe=orange!80!black,boxrule=0.6pt,arc=2pt,left=6pt,right=6pt,top=3pt,bottom=1pt,
                    title={\bfseries Je prédis, avant de commencer},coltitle=white,colbacktitle=orange!80!black,fonttitle=\small},
        retenir/.style={colback=fond,colframe=uimm,boxrule=0.6pt,arc=2pt,left=6pt,right=6pt,top=3pt,bottom=3pt,
                    title={\bfseries À retenir},coltitle=white,colbacktitle=uimm,fonttitle=\small},
        plus/.style={colback=violet!6,colframe=violet,boxrule=0.6pt,arc=2pt,left=6pt,right=6pt,top=3pt,bottom=3pt,
                    title={\bfseries Pour aller plus loin},coltitle=white,colbacktitle=violet,fonttitle=\small},
        act/.style={colback=uimm!10,colframe=uimm,boxrule=0.5pt,arc=2pt,left=5pt,right=5pt,top=2pt,bottom=2pt,before skip=7pt,after skip=3pt},
        verif/.style={colback=orange!6,colframe=orange!80!black,boxrule=0.6pt,arc=2pt,left=6pt,right=6pt,top=3pt,bottom=3pt,
                    title={\bfseries À vérifier sur poste avant la première séance},coltitle=white,colbacktitle=orange!80!black,fonttitle=\small}}
\newcommand{\fichier}[1]{\texttt{\small #1}}
\newcommand{\menu}[1]{\textsf{\textbf{#1}}}
\newcommand{\champ}[1]{\rule{#1}{0.3pt}}

\pied{Base SPC · fiche A6 · Compter les défauts : la carte c}
\begin{document}
\thispagestyle{fancy}
\entete{A6}{Compter les défauts : la carte c}{Le nombre de défauts par pièce est-il stable ?}{Cartes aux attributs}{30 min}{en binôme}{2}
\begin{tcolorbox}[boite]\textbf{Le but.} Relever des défauts par unité, calculer les limites de la carte c et la tracer.\par\textbf{Il vous faut.} La page « Tirer un lot », mode « défauts par unité » · Ellistat\end{tcolorbox}
\begin{tcolorbox}[pred]On relève en moyenne 3 défauts par panneau. Un panneau avec 8 défauts est-il anormal ?\reponse\end{tcolorbox}
\titre{À faire}
\begin{enumerate}[label=\textcolor{cyan}{\textbf{\arabic*.}},leftmargin=15pt,itemsep=1.5pt]
\item \tagW\ Page « Tirer un lot » : on compte « des défauts par unité », scénario donné par le formateur. Tirer 20 unités, une par une.
\item Noter chaque résultat et placer le point sur la carte.
\item Calculer $\bar c$ = total des défauts / nombre d'unités, puis les limites $\bar c \pm 3\sqrt{\bar c}$.
\item \tagE\ « Copier le tableau », coller dans une grille. \menu{Stats Desc.} > \menu{Cartes de contrôle aux attributs} : « Attributs » = la colonne des défauts ; laisser « Taille de lot » vide (les unités sont identiques) ; carte \textbf{c}. Dans « Phase I », Début 1 et Fin = le nombre d'unités de votre relevé.
\end{enumerate}
\vspace{3pt}{\renewcommand{\arraystretch}{1.35}\small\begin{tabularx}{\linewidth}{|l|*{10}{>{\centering\arraybackslash}X|}}\hline
\textbf{Unité} & 1 & 2 & 3 & 4 & 5 & 6 & 7 & 8 & 9 & 10\\\hline
\textbf{Défauts} &  &  &  &  &  &  &  &  &  & \\\hline\end{tabularx}}\par\vspace{2pt}
{\renewcommand{\arraystretch}{1.35}\small\begin{tabularx}{\linewidth}{|l|*{10}{>{\centering\arraybackslash}X|}}\hline
\textbf{Unité} & 11 & 12 & 13 & 14 & 15 & 16 & 17 & 18 & 19 & 20\\\hline
\textbf{Défauts} &  &  &  &  &  &  &  &  &  & \\\hline\end{tabularx}}\par\vspace{2pt}
\begin{center}\begin{tikzpicture}[x=0.780cm,y=0.31cm]\draw[gray!45,very thin] (0,0) grid (20,14);\draw[thick] (0,0)--(20,0);\draw[thick] (0,0)--(0,14);\node[left,font=\scriptsize] at (0,0) {0};\node[left,font=\scriptsize] at (0,2) {2};\node[left,font=\scriptsize] at (0,4) {4};\node[left,font=\scriptsize] at (0,6) {6};\node[left,font=\scriptsize] at (0,8) {8};\node[left,font=\scriptsize] at (0,10) {10};\node[left,font=\scriptsize] at (0,12) {12};\node[left,font=\scriptsize] at (0,14) {14};\node[below,font=\scriptsize] at (1,0) {1};\node[below,font=\scriptsize] at (2,0) {2};\node[below,font=\scriptsize] at (3,0) {3};\node[below,font=\scriptsize] at (4,0) {4};\node[below,font=\scriptsize] at (5,0) {5};\node[below,font=\scriptsize] at (6,0) {6};\node[below,font=\scriptsize] at (7,0) {7};\node[below,font=\scriptsize] at (8,0) {8};\node[below,font=\scriptsize] at (9,0) {9};\node[below,font=\scriptsize] at (10,0) {10};\node[below,font=\scriptsize] at (11,0) {11};\node[below,font=\scriptsize] at (12,0) {12};\node[below,font=\scriptsize] at (13,0) {13};\node[below,font=\scriptsize] at (14,0) {14};\node[below,font=\scriptsize] at (15,0) {15};\node[below,font=\scriptsize] at (16,0) {16};\node[below,font=\scriptsize] at (17,0) {17};\node[below,font=\scriptsize] at (18,0) {18};\node[below,font=\scriptsize] at (19,0) {19};\node[below,font=\scriptsize] at (20,0) {20};\node[rotate=90,font=\scriptsize] at (-1.41,7.0) {défauts};\node[below,font=\scriptsize] at (10.0,-1.3) {unité};\end{tikzpicture}\end{center}
\newpage
\titre{J'observe, je décide}
\begin{enumerate}[label=\textcolor{cyan}{\textbf{Q\arabic*.}},leftmargin=22pt]
\item $\bar c$ = \champ{16mm} \quad Limite supérieure = \champ{16mm} \quad Limite inférieure = \champ{16mm}\par\vspace{5pt}
\item Ellistat donne-t-il les mêmes limites ? Un point sort-il ?\reponse[1]
\item Reprenez la prédiction : 8 défauts sur un panneau, est-ce anormal ?\reponse[1]
\item Sur une carte c, que faut-il garder constant d'une unité à l'autre ?\reponse[1]
\end{enumerate}
\begin{tcolorbox}[retenir]Carte c : on compte les défauts sur des unités identiques. Les limites sont $\bar c \pm 3\sqrt{\bar c}$.\end{tcolorbox}
\begin{tcolorbox}[plus]Continuer jusqu'à 30 unités avec un scénario secret du formateur, limites figées sur les 20 premières.\end{tcolorbox}
\vfill{\footnotesize\color{uimm}\textbf{Ressources de la base :} Tirer un lot · Jeu de secours.}
\end{document}